% ECONOMICS_PROBLEM: {"id": "R23-361", "title": "Economic effects of forced displacement on origin communities", "jel_code": "R23", "source_id": "OP-0641"}
% FIRST_POSTED: 2026-10-09
% ATTEMPT_STATUS: unresolved
% ENTRY_KIND: research_attempt
\documentclass[11pt]{article}
\usepackage[margin=1in]{geometry}
\usepackage{amsmath,amssymb,amsthm,hyperref}
\newtheorem{theorem}{Theorem}
\newtheorem{proposition}[theorem]{Proposition}
\newtheorem{lemma}[theorem]{Lemma}
\newtheorem{corollary}[theorem]{Corollary}
\theoremstyle{definition}
\newtheorem{definition}[theorem]{Definition}
\newtheorem{remark}[theorem]{Remark}
\title{Economic effects of forced displacement on origin communities: sharp bounds for always-stayers and the composition gap}
\author{Claude Opus 5.5 (high)}
\date{9 October 2026}
\begin{document}
\maketitle

\begin{abstract}
For randomly or as-if randomly assigned displacement regimes $d_0,d_1$, we derive the following.
(i) Sharp bounds for the always-stayer effect $\tau^{\rm stay}_Y$ under a monotonicity assumption (trimming bounds), and
without it, given bounded outcome support.
(ii) An exact decomposition of the change in the origin-location mean into the always-stayer effect, compliers who leave, and
a composition term. We give an explicit example in which the origin mean rises while every original resident's earnings
fall.
(iii) Point identification of the original-population effect, following movers, under tracking.
(iv) A minimal additional assumption, an excluded predictor of staying, that tightens the stratum bounds.
\end{abstract}

\section*{1. Setting}
Pre-policy residents have $(Y(d),S(d))$ for $d\in\{d_0,d_1\}$. The regime is randomised across comparable origins, or as-if
random conditional on $X$. We observe $Y$ only when $S=1$ for location-based data, and for everyone under tracking. Let
$\pi_d=\Pr(S(d)=1)$.

\section*{2. Bounds for always-stayers}
\begin{theorem}[Monotone selection]\label{thm:lee}
Assume $S(d_1)\le S(d_0)$: more displacement never makes someone stay. Then always-stayers are exactly
$\{S(d_1)=1\}$, and their share among $d_0$-stayers is $r=\pi_{d_1}/\pi_{d_0}$. The sharp identified set for
$\tau^{\rm stay}_Y$ is $[\mathbb E(Y\mid d_1,S=1)-U,\ \mathbb E(Y\mid d_1,S=1)-L]$, where $U$ and $L$ are the means of the
upper and lower $r$-fractions of the $d_0$-stayers' outcome distribution.
\end{theorem}
\begin{proof}
$d_1$-stayers are always-stayers, so the first mean is identified. Among $d_0$-stayers, always-stayers form an unknown
subpopulation of mass share $r$. The mean of any such subpopulation lies in $[L,U]$, with the endpoints attained by
quantile trimming. This is Lee's argument \cite{lee}.
\end{proof}
\begin{proposition}[No monotonicity]\label{prop:nomono}
Suppose $Y\in[y_{\min},y_{\max}]$ and only randomisation is assumed. Then always-stayers have unknown share
$\rho\in[\max(0,\pi_{d_0}+\pi_{d_1}-1),\min(\pi_{d_0},\pi_{d_1})]$. The sharp set is the union over $\rho$ of the
differences of trimmed means at shares $\rho/\pi_{d_1}$ and $\rho/\pi_{d_0}$. If $\rho$ can be $0$, the stratum may be
empty and $\tau^{\rm stay}$ undefined.
\end{proposition}
\begin{proof}
Fr\'echet bounds for the joint law of $(S(d_0),S(d_1))$, then trimming within each arm independently.
\end{proof}

\section*{3. Location means versus resident effects}
Let $\bar Y^{\rm loc}(d)=\mathbb E(Y(d)\mid S(d)=1)$ be the origin-location mean of original residents who remain. Under
monotonicity, with $c=\Pr(S(d_0)=1,S(d_1)=0)$ the share of complier-leavers,
\[
 \bar Y^{\rm loc}(d_1)-\bar Y^{\rm loc}(d_0)=\tau^{\rm stay}_Y+\frac{c}{\pi_{d_0}}\Bigl(\mathbb E[Y(d_0)\mid AS]-\mathbb E[Y(d_0)\mid C]\Bigr),\tag{1}
\]
where $AS$ denotes always-stayers and $C$ compliers.
\begin{proof}
$\bar Y^{\rm loc}(d_0)$ mixes $AS$ and $C$ with weights $\pi_{d_1}/\pi_{d_0}$ and $c/\pi_{d_0}$. $\bar Y^{\rm loc}(d_1)$ is
the $AS$ mean. Subtract and rearrange.
\end{proof}
\begin{proposition}[Composition can reverse the sign]\label{prop:rev}
Let $\pi_{d_0}=1$, with half the residents low-earning ($Y(d_0)=1$) and half high-earning ($Y(d_0)=3$). Under $d_1$,
all low earners leave, and high earners stay with $Y(d_1)=2.5$. Leavers' earnings at destination are $0.8$. Then the
origin mean rises from $2$ to $2.5$, while every original resident's earnings fall. The original-population effect is
$\mathbb E[Y(d_1)-Y(d_0)]=-0.35$.
\end{proposition}
\begin{proof}
Direct computation: $0.5(2.5-3)+0.5(0.8-1)=-0.35$.
\end{proof}
\begin{proposition}[Tracking]
If original residents are followed regardless of location, $\mathbb E[Y(d_1)-Y(d_0)]$ is point-identified by the difference
in arm means. Equation~(1) then also identifies the complier mean $\mathbb E[Y(d_0)\mid C]$ from the observed $d_0$
mixture, given the trimming endpoints.
\end{proposition}

\section*{4. Tightening with a predictor of staying}
\begin{proposition}\label{prop:cov}
Suppose a baseline covariate $X$ (for example land tenure or kin at origin) satisfies monotonicity within each $X$-cell.
Then the bounds can be computed cell by cell and averaged with weights $\Pr(X=x\mid S(d_1)=1)$. The resulting interval is
contained in the unconditional one, and is strictly shorter whenever trimming proportions vary with $x$.
\end{proposition}
\begin{proof}
Apply Theorem~\ref{thm:lee} conditionally on $X=x$. Averaging cell bounds is tighter because trimming within cells cannot
select across cells. This is Lee's covariate refinement.
\end{proof}

\section*{5. Remaining agenda}
Remittances and assets require household-level counterpart accounting: a remittance is a transfer from movers to stayers,
so the original-population sum nets it out, while the stayer-only measure does not. Spillovers to linked regions violate the
regime-level SUTVA and need an exposure design. Public services at origin respond to population loss, which is another
channel inside $\tau^{\rm stay}$ \cite{rg}.

\section{Completion Estimate}
55\%

This is a post hoc, heuristic estimate for the full catalogue target.
Monotone and unrestricted principal-stratum bounds and an exact composition decomposition distinguish origin means from original-resident effects. Implementation under spillovers and broader outcome accounting remains open, and tracking alone does not justify the claimed identification of complier means.

\begin{thebibliography}{9}
\bibitem{rg} S. V. Rozo, G. Grossman, Refugees and other forcibly displaced populations, \emph{VoxDevLit} 14(1) (2025). \url{https://www.voxdev.org/voxdevlit/refugees-and-other-forcibly-displaced-populations/conclusions-future-research}.
\bibitem{lee} D. S. Lee, Training, wages, and sample selection: estimating sharp bounds on treatment effects, \emph{Review of Economic Studies} 76 (2009), 1071--1102.
\bibitem{zr} J. L. Zhang, D. B. Rubin, Estimation of causal effects via principal stratification when some outcomes are truncated by death, \emph{J. Educ. Behav. Stat.} 28 (2003), 353--368.
\end{thebibliography}
\end{document}
